\begin{tikzpicture}[
  kontrak/.style={draw,rounded corners,align=center,text width=1.85cm,
    minimum height=1.35cm,font=\scriptsize,fill=orange!12},
  hukum/.style={draw,rounded corners,align=center,text width=1.85cm,
    minimum height=1.55cm,font=\scriptsize,fill=blue!8},
  mulai/.style={draw,rounded corners,align=center,text width=1.55cm,
    minimum height=1.35cm,font=\scriptsize,fill=praditagreen!15},
  waktu/.style={align=center,font=\scriptsize\bfseries},
  jam/.style={align=center,font=\tiny},
  panah/.style={-{Stealth[length=2mm]},thick}
]
  \draw[panah] (-0.2,0) -- (11.15,0);
  \foreach \posisi in {0,2.0,4.1,6.0,8.1,10.2}
    \draw[thick] (\posisi,0.1) -- (\posisi,-0.1);
  \node[mulai] (mulai) at (0,1.15) {Kegagalan terjadi dan ditemukan};
  \node[kontrak] (a) at (2.0,1.15) {Medikara ke mitra asuransi\\ {\tiny\mdseries kontrak, misalkan}};
  \node[kontrak] (b) at (4.1,1.15) {Vendor ke Medikara\\ {\tiny\mdseries kontrak, misalkan}};
  \node[hukum] (c) at (6.0,-1.3) {Tim tanggap ke tim sektoral\\ {\tiny\mdseries Perpres 82/2022
    Pasal 13, hanya bila IIV}};
  \node[hukum] (d) at (8.1,-1.3) {Pengendali ke subjek data dan lembaga\\ {\tiny\mdseries UU PDP
    Pasal 46}};
  \node[hukum] (e) at (10.2,-1.3) {Mitra asuransi ke OJK\\ {\tiny\mdseries POJK Pasal 31, kejadian
    pada TI mitra}};
  \draw (mulai) -- (0,0.1);
  \draw (a) -- (2.0,0.1);
  \draw (b) -- (4.1,0.1);
  \draw (c) -- (6.0,-0.1);
  \draw (d) -- (8.1,-0.1);
  \draw (e) -- (10.2,-0.1);
  \node[jam] at (0,-0.42) {Jumat 20.00};
  \node[waktu] at (2.0,-0.28) {1 jam};
  \node[jam] at (2.0,-0.55) {Jumat 21.00};
  \node[waktu] at (4.1,-0.28) {12 jam};
  \node[jam] at (4.1,-0.55) {Sabtu 08.00};
  \node[waktu] at (6.0,0.4) {$1 \times 24$ jam};
  \node[jam] at (6.0,0.68) {Sabtu 20.00};
  \node[waktu] at (8.1,0.4) {$3 \times 24$ jam};
  \node[jam] at (8.1,0.68) {Senin 20.00};
  \node[waktu] at (10.2,0.4) {5 hari kerja};
  \node[jam] at (10.2,0.68) {Jumat berikutnya};
\end{tikzpicture}
